% IEEE-style skeleton for ECOS
% Compile later with a full IEEE template; this is a content scaffold.

\documentclass[conference]{IEEEtran}
\usepackage{cite}
\usepackage{amsmath,amssymb,amsfonts}
\usepackage{graphicx}
\usepackage{booktabs}
\usepackage{url}
\usepackage{hyperref}

\title{Evidence-Calibrated Open-Vocabulary Discovery of Emerging Semantic Shifts in Longitudinal Text Streams}

\author{
\IEEEauthorblockN{Anonymous Authors}
\IEEEauthorblockA{Affiliation TBD}
}

\begin{document}
\maketitle

\begin{abstract}
Lexical semantic change systems often equate embedding displacement with meaning change.
In open-vocabulary discovery, many top-ranked candidates are distributional artifacts.
We present ECOS, which treats semantic change as a falsifiable hypothesis that must survive
topic/source/frequency controls, counterfactual matching, and evidence-calibrated confidence scoring.
\end{abstract}

\begin{IEEEkeywords}
lexical semantic change, open vocabulary, counterfactual matching, calibration, longitudinal NLP
\end{IEEEkeywords}

\section{Introduction}
% Discovery gap (EMNLP 2025). Reframe: not ``how far did the embedding move?''
% but ``does the shift remain after controlling for confounders?''

\section{Related Work}
% Temporal LSC, instance-level UOT, explainable senses, AXOLOTL, discovery evaluation.

\section{Method}
\subsection{Candidate Discovery}
\subsection{Contextual and Sense Signals}
\subsection{Confounder Controls}
\subsection{Counterfactual Verification}
\subsection{ECSS and Calibration}

\section{Experiments}
\subsection{Standard LSC}
\subsection{Open-Vocabulary Discovery}
\subsection{False-Discovery Resistance}
\subsection{Calibration}
\subsection{Ablations}

\section{Results}
% Fill from data/reports/experiments/summary_compact.json

\section{Limitations}
Counterfactual matching addresses observable confounders only and does not establish causality.

\section{Conclusion}
ECOS converts open-vocabulary semantic-change detection into an evidence-calibrated discovery process.

\bibliographystyle{IEEEtran}
\begin{thebibliography}{00}
\bibitem{semeval2020} SemEval-2020 Task 1: Unsupervised Lexical Semantic Change Detection.
\end{thebibliography}

\end{document}
